\section{データ生成メカニズムとモデリング}
\subsection{授業内容}
\subsection{授業内容}
\subsubsection{科目の中でのこのコマの位置づけ}

プログラミングの本質は代入，反復，条件分岐である。この3点について理解し，基本的な書き方を学ぶ。実際にRでコードを書きながら，その挙動について確認する。
最終的な到達段階として，Fizz-Buzz問題や行列計算ができるコードを書くこととする。

\subsubsection{コマ主題細目}
\begin{description}
	\item[高級言語の基本的]高級言語と呼ばれるプログラミング言語の基本的な働きは，代入，反復，条件分岐である。Rでは\verb|<-|や\verb|=|で代入を，\verb|for|や\verb|while|で反復を，\verb|if|や\verb|if_else|で条件分岐を行う。これらの表現は言語間を通じて共通なものが多いため，その基本的な振る舞いを確認することは技術の一般化に役立つ。
	\item[オブジェクトの型] Rは変数を宣言せずに利用できるところが初心者向けの長所ではあるが，Rの中にストックされるオブジェクトの種類や型について正しく理解しておくことが，今後のデバッグの上で役にたつ。変数としての数値型(Int,Double,Complex)，文字型(Char,Factor)，論理型(logical)や，オブジェクトのベクトル，行列，配列，リスト，データフレームについて理解する。
	\item[関数を作る] コマンド(文)は連なってスクリプトとなる。反復されるスクリプトは関数化すると効率的である。ここではRを用いた関数の定義の仕方について学ぶ。引数やデフォルトの値について知ることは，ヘルプを見る時の役にも立つ。また関数化はスパゲティ・コードになることを避ける，カプセル化された表記方法への第一歩であり，オブジェクト指向プログラミングへの第一歩でもある。
\end{description}
\subsubsection{キーワード}
\begin{itemize}
	\item オブジェクトの型
	\item 代入
	\item 反復
	\item 条件分岐
	\item 関数
\end{itemize}
\subsection{授業情報}
\paragraph{コマの展開方法}
講義/遠隔/演習
\subsubsection{予習・復習課題}
\paragraph{予習}RStudioにおけるプロジェクト単位での管理など，R/RStudioの基本的な使い方については以後も逐一支持することはないので，これまでのことをしっかり復習し，本講義の予習とする。
\paragraph{復習}反復計算の練習課題，条件分岐の練習課題など，複数の課題にしっかり取り組むこと。